\documentclass{subfiles}
\begin{document}\label{Sec:1.2}
    As stated previously, in some cases the worst cases analysis is missleading.
        Consider \Cref{Algo:1}, it is obvious that this has a worst case time complexity of \(\bigO{n}\).
        What we are interested in is the following: given an algorithm, what is its real time complexity 
        over a sequence of \(k\) iterations?
        \subfile{../../TikZ/Algorithm - Binary counter}

    To answer this question, we have to define a cost function that depends on \(k\).
        Let 
        \[
            T_{\alpha}(n) = \frac{T(n, k)}{k}
        \]
        be the \emph{amortized cost} of a single iteration, 
        where \(T(n, k)\) is the maximum of the worst cases among all \(k\) iterations.

    In the case of our binary counter example, we have the following: 
        the first bit is flipped \(k\) times, the second one half of that, 
        the third a half of the second, etc. Summing up, we have that 
        \[\begin{aligned}
             T(n, k) & = k + \floor*{\frac{k}{2}} + \floor*{\frac{k}{4}} + \cdots + 2 + 1 \\
                & = \sum_{i = 0}^{\floor*{\log_{2} k}} \le k \sum_{i = 0}^{\infty} 2^{-i} = 2k.
        \end{aligned}\] 
        Hence, the amortized time is \(T_{\alpha}(n) = \sfrac{2k}{k} = 2 = \bigO{1}\),
            while the worst case was \(\bigO{n}\).

    It should be clear by now what amortized analysis is about. If not, 
        we can summarize it as follow: amortized analysis allows us to show that, 
        over a sequence of operations, the cost of a single operation is much lower 
        then what the worst case analysis suggest.


    In the next two sections we provide two methods to compute the amortized times.

    \subsubsection{Accounting method}
    \subfile{../subsub/1.2.1 - Accounting method}

    \subsubsection{Potential method}
    \subfile{../subsub/1.2.2 - Potential method}
    \clearpage
\end{document}
